\chapter{Installation}
\label{chap:installation}

This chapter explains how to build and install Medorama from source, on Linux, macOS and
Windows, with or without administrator rights or internet access. It is written for whoever
installs Medorama: a researcher on a workstation, or an administrator on a server. The quickest
route is one of the installers (\cref{sec:installation-installers}); the rest of the chapter
builds each program by hand. Once Medorama is built, \cref{chap:quickstart} sets up a catalog
with the demo, and \cref{chap:administration} sets one up for real data. Installing the viewers
on their own, without Medorama, is described in \cref{chap:viewers-alone}. Copy-and-paste
commands for each operating system, for Medorama and for each viewer on its own, are also on
the website, at \url{https://medorama.org/docs/install.html} (\file{docs/install.md}).

\section{What there is to build}
\label{sec:installation-programs}

\Cref{tab:installation-programs} lists the programs of the repository.

\begin{table}[htbp]
\caption{The programs, and what each is built from.}
\label{tab:installation-programs}
\begin{tabularx}{\linewidth}{@{}P{2.6cm}P{1.4cm}P{4.2cm}L@{}}
\toprule
{\tabhead Program} & {\tabhead Language} & {\tabhead Built from} & {\tabhead Result} \\
\midrule
\cmd{medorama} & Rust & The workspace (\cmd{cargo build -p medorama}) & One binary (about 18\,MB)
  with the database engine, the web page and the ECG and PACS viewers \\
\cmd{vitalstream-rs} & Rust & The workspace (\opt{-p vitalstream-rs}) & One binary with the web
  page \\
\cmd{vitalstream-cpp} & C++20 & \file{viewers/vitalstream/cpp} (CMake) & One binary with the web
  page \\
\cmd{vitalstream-py} & Python & \file{viewers/vitalstream/python/vitalstream} & A Python
  package \\
\cmd{muse-py}, \cmd{dicom-py} & Python & \file{viewers/ecg/python/muse},
  \file{viewers/pacs/python/dicom} & The database tools: Python packages \\
\bottomrule
\end{tabularx}
\end{table}

WFDB records are read with wfdb-x in each language: Cargo, CMake and pip fetch it from GitHub at
one pinned commit (\code{a707ba9}), like any other dependency. Every program can also be built
without WFDB support, and then contains no wfdb-x code (\cref{sec:installation-no-wfdb}). The
ECG and PACS viewers are static pages: there is nothing to build. Add-ons (adapters) are built in their own repositories
(\cref{chap:addons}).

Every Rust program is a member of one Cargo workspace (\file{Cargo.toml} at the root of the
repository), sharing one \file{Cargo.lock} and one \file{target/} folder. Each builds and
deploys on its own.

\section{Prerequisites}
\label{sec:installation-prerequisites}

\Cref{tab:installation-prerequisites} lists what each system needs. Only the compiler may need
an administrator; Rust, CMake, Ninja, the Python packages and every program built here install
in the user's account (\cref{sec:installation-no-admin}).

\begin{table}[htbp]
\caption{Prerequisites, per system.}
\label{tab:installation-prerequisites}
\begin{tabularx}{\linewidth}{@{}P{2.4cm}LP{3.0cm}L@{}}
\toprule
& {\tabhead Linux} & {\tabhead macOS} & {\tabhead Windows 10/11} \\
\midrule
Git, Python 3.9+ & \cmd{sudo apt install git python3 python3-venv} (Debian, Ubuntu) or
  \cmd{sudo dnf install git python3} (Fedora, RHEL) & \cmd{xcode-select --install} (the
  Command Line Tools: Git, Python, compilers) & \cmd{winget install Git.Git} and
  \cmd{winget install Python.Python.3.12} \\
A C compiler, for Rust (SQLite is compiled in) & \cmd{sudo apt install build-essential} or
  \cmd{sudo dnf install gcc} & The Command Line Tools & The Microsoft C++ Build Tools (below) \\
A C++20 compiler, for the C++ viewer & GCC 12+ or Clang 15+ (\cmd{sudo apt install g++},
  Debian 12 and Ubuntu 22.04 or later) & Apple Clang 15+ (the Command Line Tools) & Visual
  Studio 2022 (the Build Tools) \\
Rust 1.88+ & rustup (\url{https://rustup.rs}), in the user's account; the installers install
  it when missing & The same & The same \\
CMake 3.24+, Ninja (C++) & \cmd{pip install "cmake>=3.24" ninja}; the installers do it, into
  \file{.venv} & The same & The same \\
\bottomrule
\end{tabularx}
\end{table}

On Windows, the Microsoft C++ Build Tools are installed with:
\begin{lstlisting}[language=PowerShell]
winget install Microsoft.VisualStudio.2022.BuildTools `
  --override "--passive --wait --add Microsoft.VisualStudio.Workload.VCTools --includeRecommended"
\end{lstlisting}
Their workload \emph{Desktop development with C++} is what both Rust and C++ need.

\section{The installers}
\label{sec:installation-installers}

Three installers build what is needed into the clone. They need no administrator rights,
except, on some systems, for the compiler (\cref{sec:installation-prerequisites}). Each has a
shell script for Linux and macOS and a PowerShell script for Windows
(\cref{tab:installation-installers}).

\begin{table}[htbp]
\caption{The installers.}
\label{tab:installation-installers}
\begin{tabularx}{\linewidth}{@{}P{5.2cm}L@{}}
\toprule
{\tabhead Installer} & {\tabhead Builds} \\
\midrule
\file{tools/install-medorama.sh} (Windows: \file{tools\textbackslash install-medorama.ps1}) &
  Medorama, \emph{one} VitalStream viewer (the C++ viewer by default) and the database tools,
  into \file{.venv/}; the viewer is the one Medorama runs. The installer for using Medorama. \\
\file{viewers/vitalstream/install.sh} (Windows: \file{viewers\textbackslash vitalstream\textbackslash install.ps1})
  & VitalStream alone, without Medorama, into \file{\textasciitilde/.local/bin} (Windows:
  \file{\%LOCALAPPDATA\%\textbackslash vitalstream\textbackslash bin}), with a root folder it
  never leaves (by default the home folder): \cref{chap:viewers-alone}. \\
\file{tools/install.sh} (Windows: \file{tools\textbackslash install.ps1}) & Everything, for
  development: the Python packages, the Rust workspace, the C++ viewer and its tests, and the
  repository's git hooks. \\
\bottomrule
\end{tabularx}
\end{table}

Run them from the root of a clone. On Linux or macOS:
\begin{lstlisting}[style=shell]
git clone https://github.com/alphanumericslab/medorama.git && cd medorama
tools/install-medorama.sh
source .venv/bin/activate       # in every new terminal
medorama --version
\end{lstlisting}
On Windows, in PowerShell:
\begin{lstlisting}[language=PowerShell]
git clone https://github.com/alphanumericslab/medorama.git; cd medorama
powershell -ExecutionPolicy Bypass -File tools\install-medorama.ps1
.venv\Scripts\Activate.ps1
medorama --version
\end{lstlisting}

The programs go into \file{.venv/bin} (Windows: \file{.venv\textbackslash Scripts}), which is
not on the \code{PATH} by itself: activate the environment in each new terminal, or call
\file{.venv/bin/medorama} by its path.

\subsection{The options}
\label{sec:installation-options}

\Cref{tab:installation-options} lists the options of the three installers. The PowerShell
scripts take the same options, written as \opt{-Viewer}, \opt{-NoTools}, \opt{-Demo} and so on.

\begin{xltabular}{\linewidth}{@{}P{3.4cm}P{2.6cm}L@{}}
\caption{The installers' options.}
\label{tab:installation-options}\\
\toprule
{\tabhead Option} & {\tabhead Installers} & {\tabhead Effect} \\
\midrule
\endfirsthead
\toprule
{\tabhead Option} & {\tabhead Installers} & {\tabhead Effect} \\
\midrule
\endhead
\bottomrule
\endfoot
\opt{--viewer cpp|rs|py|none} (\opt{-Viewer}) & \file{install-medorama} & The viewer Medorama runs:
  \code{cpp} (the default) needs a C++20 compiler; \code{rs} needs only Rust; \code{py} the
  Python viewer and its packages; \code{none}, Medorama alone. The choice is recorded in
  \file{.venv/bin/medorama.viewer} (\cref{chap:viewers-in-medorama}). \\
\opt{--no-tools} (\opt{-NoTools}) & \file{install-medorama} & Leave out the database tools
  (\cmd{muse-py}, \cmd{dicom-py}). \\
\opt{--viewer cpp|rs|py} (\opt{-Viewer}) & VitalStream's \file{install} & The viewer to build (default
  \code{cpp}). \\
\opt{--prefix DIR} (\opt{-Prefix}) & VitalStream's \file{install} & Where the viewer goes:
  \file{DIR/bin} (default \file{\textasciitilde/.local}; Windows
  \file{\%LOCALAPPDATA\%\textbackslash vitalstream}). The Python viewer's environment goes in
  \file{DIR/share/vitalstream/venv} (Windows \file{DIR\textbackslash venv}). \\
\opt{--root DIR} (\opt{-Root}) & VitalStream's \file{install} & The viewer's root, fixed in the
  program: it never opens anything outside it. Default \code{\textasciitilde}, the home folder
  of whoever runs it; \code{/} for the whole file system. Medorama's installer gives its viewer
  the whole file system: Medorama's own rules decide what each user opens. \\
\opt{--no-password} (\opt{-NoPassword}) & VitalStream's \file{install} & Build the viewer
  without password support: no password can be set, and it refuses to run where one is
  required (under Medorama). Default: password support, a password being optional. \\
\opt{--py}, \opt{--rs}, \opt{--cpp} (\opt{-Py}, \opt{-Rs}, \opt{-Cpp}) & \file{install} &
  Only these parts. The Python packages are always installed: \opt{--py} alone installs only
  them, \opt{--rs} adds the Rust workspace, \opt{--cpp} the C++ viewer, \opt{--rs --cpp} both.
  Without any of the three: all three. \\
\opt{--no-rustup} (\opt{-NoRustup}) & \file{install-medorama.sh} and \file{.ps1},
  \file{install.sh} and \file{.ps1}, VitalStream's \file{install.sh} & Fail rather than install
  Rust when \cmd{cargo} is missing. VitalStream's \file{install.ps1} has no \opt{-NoRustup}: it
  never installs Rust, and \opt{-Viewer rs} fails when \cmd{cargo} is missing. \\
\opt{--demo} (\opt{-Demo}) & \file{install-medorama}, \file{install} & Also make the demo:
  invented patients that \cmd{medorama init} loads into a cohort \code{DEMO}
  (\cref{chap:quickstart}). Off by default. \\
\opt{--no-wfdb} (\opt{-NoWfdb}) & \file{install-medorama}, VitalStream's \file{install} & Build
  without WFDB support, with no wfdb-x code: a WFDB recording is refused, every other format is
  read (\cref{sec:installation-no-wfdb}). \file{install-medorama} builds both Medorama and its
  viewer this way, and refuses \opt{--demo} with it (the demo's recordings are WFDB records). \\
\opt{--help} & The shell scripts & Print the script's description and options. \\
\end{xltabular}

All three use the compilers named in \code{CC} and \code{CXX} when they are set (after
\cmd{module load gcc}, or in a conda environment), so that no administrator rights are needed
when such a compiler is available (\cref{sec:installation-no-admin}). On macOS, when Xcode is
installed but its license was never accepted, they build with the Command Line Tools instead.

\subsection{What the Medorama installer does}
\label{sec:installation-what}

\file{tools/install-medorama.sh} takes these steps, in order:
\begin{enumerate}
\item It makes the Python environment \file{.venv} (with Python 3.9 or later) if it is not
  there.
\item It installs Rust in the user's account (\file{\textasciitilde/.rustup},
  \file{\textasciitilde/.cargo}) if \cmd{cargo} is missing, unless \opt{--no-rustup} is given.
\item It builds Medorama (\cmd{cargo build --release --locked -p medorama}) and links it into
  \file{.venv/bin}.
\item It builds the chosen viewer, with the whole file system as its root and with password
  support, and links it into \file{.venv/bin}. With \opt{--no-wfdb}, Medorama and the viewer are
  built without WFDB support (\cref{sec:installation-no-wfdb}). The C++ viewer is built in
  \file{viewers/vitalstream/cpp/build/viewer}, after CMake and Ninja are installed into
  \file{.venv} (a build folder configured with another compiler is removed and made again); the
  Rust viewer is built in \file{target/release}; the Python viewer is installed as an editable
  package in \file{.venv}.
\item It installs the database tools \cmd{muse-py} and \cmd{dicom-py} into \file{.venv}, unless
  \opt{--no-tools} is given.
\item With \opt{--demo}, it makes the demo (\cref{sec:installation-demo}). Without it, it removes
  \file{.venv/bin/medorama.demo}, so that catalogs created afterwards start empty.
\item It records the viewer chosen in \file{.venv/bin/medorama.viewer}.
\item In a git clone, it enables the repository's git hooks, which run the privacy check
  (\file{tools/git-hooks/check-files}) before each commit and merge; it warns if the clone
  already uses other hooks.
\end{enumerate}
Cargo builds into the clone's \file{target/}, whatever \code{CARGO\_TARGET\_DIR} or a
\code{build.target-dir} says, and the installer stops with an error if a program it installs is
missing at the end.
It ends by printing the commands it installed and the next steps:
\cmd{source .venv/bin/activate}, \cmd{medorama init --results DIR} and \cmd{medorama serve}.

\subsection{The demo}
\label{sec:installation-demo}

The \emph{demo} is a set of six invented patients with their encounters, laboratory results,
vital signs, medications, notes, ICU recordings, a 24-hour Holter, an overnight sleep study (an
EDF+ file), ECGs and imaging (\cref{chap:quickstart}). With \opt{--demo}, the Medorama and
development installers also take these steps:
\begin{enumerate}
\item They install NumPy and wfdb-x into \file{.venv} if they are not there: the demo's WFDB
  records are written with wfdb-x, at the commit that the Python viewer uses.
\item They write the demo's files into \file{.venv/share/medorama/demo} (Windows
  \file{.venv\textbackslash share\textbackslash medorama\textbackslash demo}; about 60\,MB),
  with \file{tools/make-demo-data.py}.
\item They name that folder in \file{.venv/bin/medorama.demo} (Windows
  \file{.venv\textbackslash Scripts\textbackslash medorama.demo}), next to the programs.
\end{enumerate}
\cmd{medorama init} then loads the demo by itself (\cref{chap:quickstart}). Running the
installer again without \opt{--demo} removes \file{medorama.demo}: catalogs created after that
start empty, and a catalog that holds the demo already keeps it. The demo needs WFDB support:
its ICU and Holter recordings are WFDB records, so \opt{--demo} and \opt{--no-wfdb} are refused
together.

\begin{note}
The demo is for trying, demonstrating and teaching Medorama. Leave it out of a server meant for
real patients (\cref{chap:deployment}).
\end{note}

\section{Building by hand}
\label{sec:installation-by-hand}

\subsection{The Rust workspace}
\label{sec:installation-rust}

\paragraph{Rust,} in the user's account (open a new shell afterwards):
\begin{lstlisting}[style=shell]
curl --proto '=https' --tlsv1.2 -sSf https://sh.rustup.rs | sh -s -- -y --profile minimal
rustc --version                          # 1.88 or later
\end{lstlisting}
On Windows: \cmd{winget install Rustlang.Rustup}. Rust 1.88 is the oldest version that builds
the workspace with the versions of its \file{Cargo.lock} (\code{rust-version} in
\file{Cargo.toml}; CI checks it, and \cmd{make msrv} does locally). Development and CI use the
current stable Rust; the released programs are built with a pinned version.

\paragraph{Build,} from the root of the clone:
\begin{lstlisting}[style=shell]
cargo build --release --locked                  # medorama, vitalstream-rs -> target/release/
cargo build --release --locked -p medorama      # one program
cargo install --path medorama --locked          # optional: into ~/.cargo/bin (on the PATH)
cargo install --path viewers/vitalstream/rust --locked
\end{lstlisting}

\paragraph{For an older Linux server.} A binary built on a recent distribution needs a C library
at least as recent. A fully static build runs on any Linux of the same architecture:
\begin{lstlisting}[style=shell]
rustup target add x86_64-unknown-linux-musl
sudo apt install musl-tools                     # Fedora: sudo dnf install musl-gcc
cargo build --release --locked --target x86_64-unknown-linux-musl -p medorama
# -> target/x86_64-unknown-linux-musl/release/medorama
\end{lstlisting}

\paragraph{Database encryption} (SQLCipher, with OpenSSL built from source; needs Perl and
\cmd{make}):
\begin{lstlisting}[style=shell]
cargo build --release --locked -p medorama --features encryption
\end{lstlisting}
The key and the file that holds it are described in \cref{chap:deployment}.

\paragraph{macOS.} If the build stops with \code{linking with cc failed: exit status: 69} or
``You have not agreed to the Xcode license agreements'', Xcode is installed but its license was
never accepted: run \cmd{sudo xcodebuild -license accept}, or build with the Command Line Tools:
\begin{lstlisting}[style=shell]
DEVELOPER_DIR=/Library/Developer/CommandLineTools cargo build --release --locked
\end{lstlisting}

\paragraph{Windows} builds in an ordinary PowerShell window once the Build Tools are installed;
the programs are \file{target\textbackslash release\textbackslash *.exe}.

\subsection{The C++ viewer}
\label{sec:installation-cpp}

From \file{viewers/vitalstream/cpp}, with a C++20 compiler, CMake 3.24+ and Ninja:
\begin{lstlisting}[style=shell]
cmake --preset release && cmake --build --preset release  # -> build/release/bin/vitalstream-cpp
ctest --preset release                                    # the tests
cmake --install build/release --prefix ~/.local           # optional: ~/.local/bin/vitalstream-cpp
\end{lstlisting}

\begin{table}[htbp]
\caption{The C++ viewer's presets and options.}
\label{tab:installation-cpp}
\begin{tabularx}{\linewidth}{@{}P{3.6cm}P{2.6cm}L@{}}
\toprule
{\tabhead Preset} & {\tabhead Folder} & {\tabhead Builds} \\
\midrule
\code{release} & \file{build/release} & The viewer and its tests \\
\code{debug} & \file{build/debug} & The same, unoptimized, with debug information \\
\code{viewer} & \file{build/viewer} & The viewer only (what the installers use) \\
\code{nowfdb} & \file{build/nowfdb} & The viewer and its tests without WFDB support
  (\code{VS\_WFDB=OFF}; \cref{sec:installation-no-wfdb}) \\
\midrule
{\tabhead Option (\code{-D\dots})} & {\tabhead Default} & {\tabhead Effect} \\
\midrule
\code{VS\_BUILD\_VIEWER} & \code{ON} & The viewer (off: only the libraries \code{vs\_core} and
  \code{vs\_engine}, for adapters written in C++) \\
\code{VS\_BUILD\_TESTS} & \code{ON} & The unit tests (GoogleTest) \\
\code{VS\_DEPS\_DIR} & \file{cpp/deps} & Pre-downloaded dependency archives
  (\cref{sec:installation-offline}) \\
\code{VS\_ROOT} & (unset) & The viewer's root (\cref{sec:installation-viewer-options}) \\
\code{VS\_PASSWORD} & \code{ON} & Password support \\
\code{VS\_WFDB} & \code{ON} & WFDB support (off: wfdb-x neither downloaded nor built in;
  \cref{sec:installation-no-wfdb}) \\
\bottomrule
\end{tabularx}
\end{table}

The first configuration downloads five pinned archives (wfdb-x, CLI11, toml++, cpp-httplib;
GoogleTest for the tests; four with \code{VS\_WFDB=OFF}), each checked against its SHA-256
(\file{viewers/vitalstream/cpp/cmake/deps.cmake}). Without Ninja:
\begin{lstlisting}[style=shell]
cmake -S . -B build/make -G "Unix Makefiles" -DCMAKE_BUILD_TYPE=Release
cmake --build build/make
\end{lstlisting}
and on Windows \cmd{cmake -S . -B build/vs -G "Visual Studio 17 2022"} then
\cmd{cmake --build build/vs --config Release}.

\paragraph{Linux.} \cmd{g++ --version} must say 12 or later. On older systems, install a newer
GCC (\cmd{sudo apt install g++-12}; RHEL 8: \cmd{sudo dnf install gcc-toolset-12}) and select it
with \cmd{export CC=gcc-12 CXX=g++-12} (RHEL 8: \cmd{scl enable gcc-toolset-12 bash}), or use
Clang 15+ (\cmd{export CC=clang-15 CXX=clang++-15}). Without administrator rights, follow
\cref{sec:installation-no-admin}. To run a program built against a newer C++ runtime on other
machines, link the runtime in:
\code{-DCMAKE\_EXE\_LINKER\_FLAGS="-static-libstdc++ -static-libgcc"}.

\paragraph{macOS.} Apple Clang 15+ comes with the Command Line Tools
(\cmd{xcode-select --install}, once, with an administrator's password). With GCC from Homebrew
(\code{CC=gcc-14 CXX=g++-14}), add \code{-DVS\_BUILD\_TESTS=OFF}: GoogleTest does not compile
with GCC against the macOS 26 SDK. If a standard header such as \code{cstddef} is not found, the
Command Line Tools are damaged: reinstall them with
\cmd{sudo rm -rf /Library/Developer/CommandLineTools \&\& xcode-select --install} (meanwhile, the
build detects this case and uses the SDK's headers).

\paragraph{Windows.} In a \emph{Developer PowerShell for VS 2022} (Start menu), which sets up the
compiler:
\begin{lstlisting}[language=PowerShell]
cd viewers\vitalstream\cpp
cmake --preset release; cmake --build --preset release  # -> build\release\bin\vitalstream-cpp.exe
\end{lstlisting}
In an ordinary PowerShell window, \cmd{cmake --preset release} stops with ``No
CMAKE\_CXX\_COMPILER could be found''; \cmd{tools\textbackslash install.ps1 -Cpp} sets the
compiler up itself. Visual Studio can also open the \file{cpp} folder directly (it reads
\file{CMakePresets.json}).

\subsection{The Python packages}
\label{sec:installation-python}

\begin{lstlisting}[style=shell]
python3 -m venv .venv
.venv/bin/pip install './viewers/vitalstream/python/vitalstream[wfdb]' # vitalstream-py
.venv/bin/pip install ./viewers/ecg/python/muse ./viewers/pacs/python/dicom   # muse-py, dicom-py
\end{lstlisting}
The \code{[wfdb]} extra installs wfdb-x from GitHub (the archive of the pinned commit; no
\cmd{git} needed); without it, \cmd{vitalstream-py} has no WFDB support
(\cref{sec:installation-no-wfdb}), and the installers always add it, except with
\opt{--no-wfdb}. \cmd{muse-py} and \cmd{dicom-py} also serve the ECG and PACS viewers
(\cmd{serve}: \cref{chap:viewers-alone}). For development, install them editable
(\cmd{pip install -e \dots{} --config-settings editable\_mode=compat}) with their
\code{[test]} extras (\code{[test,wfdb]} for \cmd{vitalstream-py}), as \file{tools/install.sh}
does.

Each package carries the license: building it copies the repository's \file{LICENSE} (with its
Required Notice) and \file{COMMERCIAL.md} beside its \file{pyproject.toml}, and they are
installed with the package (in its \file{.dist-info/licenses/}). A package builds from a clone,
or from its source distribution: \cmd{python -m build viewers/vitalstream/python/vitalstream}
(and \file{viewers/ecg/python/muse}, \file{viewers/pacs/python/dicom}) writes an sdist that
carries the viewer's page and the license, and a wheel built from that sdist alone. Building
needs setuptools 77 or later, which pip fetches into the isolated build environment.

\subsection{A viewer's build options}
\label{sec:installation-viewer-options}

A viewer installed on its own has two options fixed when it is built: its root, and whether it
supports a password (WFDB support is a third, for every program: \cref{sec:installation-no-wfdb}).
The VitalStream installer passes them (\opt{--root}, \opt{--no-password});
a build by hand sets them as \cref{tab:installation-viewer-options} shows. A root may be written
with \code{\textasciitilde} (the home folder of whoever runs the viewer); unset, the root is the
whole file system.

\begin{table}[htbp]
\caption{A viewer's root and password support, set when it is built.}
\label{tab:installation-viewer-options}
\begin{tabularx}{\linewidth}{@{}P{2.8cm}LP{3.4cm}@{}}
\toprule
{\tabhead Viewer} & {\tabhead Root} & {\tabhead Without password support} \\
\midrule
\cmd{vitalstream-rs} & \code{VITALSTREAM\_ROOT=DIR} \cmd{cargo build --release -p vitalstream-rs
  --target-dir target/standalone} &
  \opt{--no-default-features --features wfdb} (the Cargo feature \code{password} left out) \\
\cmd{vitalstream-cpp} & \code{cmake --preset viewer -B build/standalone -DVS\_ROOT=DIR} &
  \code{-DVS\_PASSWORD=OFF} \\
\cmd{vitalstream-py} & \code{VITALSTREAM\_ROOT=DIR} \cmd{pip install}
  \file{'./viewers/vitalstream/python/vitalstream[wfdb]'} & \code{VITALSTREAM\_PASSWORD=0} \\
\cmd{muse-py serve}, \cmd{dicom-py serve} (the ECG and PACS viewers served) &
  \code{VIEWER\_ROOT=DIR} \cmd{pip install} \file{./viewers/ecg/python/muse} (or
  \file{\dots/pacs/python/dicom}) & \code{VIEWER\_PASSWORD=0} \\
\bottomrule
\end{tabularx}
\end{table}

At run time, the settings (\code{data\_roots}, \opt{--root}) narrow the root further, never
widen it. The standalone installer builds in folders of its own (\file{target/standalone},
\file{viewers/vitalstream/cpp/build/standalone}), so that Medorama's viewer
(\file{target/release}, \file{build/viewer}) keeps its own options.

\section{Without WFDB}
\label{sec:installation-no-wfdb}

WFDB support is the only part of the programs that uses wfdb-x. Each program can be built
without it, and such a build contains no wfdb-x code. It reads every other format as usual: EDF,
EDF+, BDF and BDF+ files, the formats of installed add-ons (such as STP, GE MUSE and DICOM), CLIF
tables, and the databases Medorama imports or connects to. A WFDB recording (a \file{.hea}
file, a multi-segment record, a folder of records or a WFDB collection) is still recognised, and
refused with one message, never skipped in silence:
\begin{lstlisting}[style=text]
this build has no WFDB support (built without the `wfdb` feature)
\end{lstlisting}
The C++ viewer ends the message with ``(built with VS\_WFDB=OFF)'', and the Python viewer with
``(wfdb-x is not installed: pip install 'vitalstream\_py[wfdb]')''. Medorama does not add such a
recording: an import lists it among the files it skipped, with this reason. A viewer lists it
when browsing, and refuses to open or scan it (HTTP 422, with the message). At start-up, every
program built this way says so. \Cref{tab:installation-no-wfdb} shows how each program is built
without WFDB support.

\begin{table}[htbp]
\caption{Building without WFDB support.}
\label{tab:installation-no-wfdb}
\begin{tabularx}{\linewidth}{@{}P{2.8cm}LP{3.0cm}@{}}
\toprule
{\tabhead Program} & {\tabhead Without WFDB support} & {\tabhead The default} \\
\midrule
\cmd{medorama} & \cmd{cargo build --release --locked -p medorama --no-default-features} (add
  \opt{--features encryption} for database encryption) & The Cargo feature \code{wfdb}, on \\
\cmd{vitalstream-rs} & \cmd{cargo build --release --locked -p vitalstream-rs
  --no-default-features --features password} & The Cargo features \code{password} and
  \code{wfdb}, on \\
\code{recording-engine} (for adapters written in Rust) & \code{default-features = false} in the
  adapter's \file{Cargo.toml} & The Cargo feature \code{wfdb}, on \\
\cmd{vitalstream-cpp} & From \file{viewers/vitalstream/cpp}: \cmd{cmake --preset nowfdb \&\&
  cmake --build --preset nowfdb} (into \file{build/nowfdb}), or \code{-DVS\_WFDB=OFF} with any
  preset & \code{VS\_WFDB=ON} \\
\cmd{vitalstream-py} & \cmd{pip install ./viewers/vitalstream/python/vitalstream}, without the
  \code{[wfdb]} extra, in an environment that does not hold wfdb-x (\cmd{pip uninstall wfdb-x}
  removes it) & The installers install the \code{[wfdb]} extra \\
\bottomrule
\end{tabularx}
\end{table}

The installers do the same with one option, \opt{--no-wfdb} (Windows: \opt{-NoWfdb}):
\file{tools/install-medorama.sh --no-wfdb} builds both Medorama and its viewer this way, and
\file{viewers/vitalstream/install.sh --no-wfdb} one viewer alone. Medorama and its viewer are
built separately, so a deployment without wfdb-x builds both without WFDB support. The demo
needs WFDB support, because its recordings are WFDB records written with wfdb-x.

To check a build, \cmd{cargo tree --locked -p medorama --no-default-features | grep wfdbx}
prints nothing (likewise for \cmd{vitalstream-rs}), a C++ build in \file{build/nowfdb} has no
\file{\_deps/wfdbx-src} folder, and \cmd{make test-nowfdb} runs the tests of these builds (the
Python tests with wfdb-x made unimportable). The tests on WFDB records are skipped there, and a
test of each program checks the refusal.

Two parts read WFDB headers without wfdb-x, so they work in either build: the ECG viewer, which
shows a WFDB record attached to an item with its own reader in the page, and Privacy Mode, which
de-identifies the text of a WFDB header. In the C++ libraries, \file{core/wfdb\_c.hpp} (WFDB
headers, records and annotations, for adapters written in C++) exists only with
\code{VS\_WFDB=ON}.

\section{Without administrator rights}
\label{sec:installation-no-admin}

Only the compiler may need an administrator. Rust, CMake, Ninja, the Python packages and every
program built here install in the user's account.

\paragraph{Linux (a shared server, a cluster).} \cmd{g++ --version} 12 or later is enough.
Otherwise, use the first of the ways in \cref{tab:installation-gcc} that the machine offers.

\begin{table}[htbp]
\caption{A C++20 compiler without administrator rights, on Linux.}
\label{tab:installation-gcc}
\begin{tabularx}{\linewidth}{@{}P{3.2cm}LL@{}}
\toprule
{\tabhead Way} & {\tabhead Check} & {\tabhead If available} \\
\midrule
A compiler module & \cmd{module avail 2>\&1 | grep -i gcc} lists GCC 12 or later &
  \cmd{module load <that name>}, then \cmd{export CC=gcc CXX=g++} \\
The gcc-toolset (RHEL, Rocky, Alma 8 and 9) & \cmd{ls /opt/rh} lists \code{gcc-toolset-12} or later &
  \cmd{scl enable gcc-toolset-13 bash} \\
GCC 13 in the home folder (anywhere) & & \cmd{tools/install-gcc.sh}, then
  \cmd{source \textasciitilde/miniforge3/bin/activate vs-build} in every new shell \\
\bottomrule
\end{tabularx}
\end{table}

\file{tools/install-gcc.sh} installs Miniforge (conda-forge) into
\file{\textasciitilde/miniforge3}, with an environment \code{vs-build} that holds GCC 13, CMake,
Ninja and Python 3.12 (about 1\,GB). The installers use the compiler in \code{CC} and \code{CXX}
when they are set. These hold for the current shell only: repeat the \cmd{module load},
\cmd{scl} or \cmd{source} line in every new shell (or add it to \file{\textasciitilde/.bashrc}).
CMake remembers the compiler a build folder was first configured with: after switching
compilers, delete the build folder (the installers do it).

\paragraph{CMake and Ninja} in the user's account:
\cmd{python3 -m pip install --user "cmake>=3.24" ninja}, or the official archives
(\url{https://github.com/Kitware/CMake/releases},
\url{https://github.com/ninja-build/ninja/releases}) unpacked into
\file{\textasciitilde/.local}.

\paragraph{Windows.} For Rust, the self-contained GNU toolchain needs no Visual Studio:
\cmd{rustup-init.exe --default-host x86\_64-pc-windows-gnu} (SQLite then needs GCC from MSYS2).
For C++, MSYS2 (\url{https://www.msys2.org}) unpacked into the user's folder provides GCC,
CMake and Ninja:
\begin{lstlisting}[language=PowerShell]
$base = "https://github.com/msys2/msys2-installer/releases/download/nightly-x86_64"
Invoke-WebRequest "$base/msys2-base-x86_64-latest.sfx.exe" -OutFile msys2.exe
.\msys2.exe -y -o"$env:USERPROFILE"; Remove-Item msys2.exe
$bash = "$env:USERPROFILE\msys64\usr\bin\bash.exe"
& $bash -lc "pacman -Syu --noconfirm"
$pkgs = "git mingw-w64-ucrt-x86_64-gcc mingw-w64-ucrt-x86_64-cmake mingw-w64-ucrt-x86_64-ninja"
& $bash -lc "pacman -S --noconfirm --needed $pkgs"
\end{lstlisting}
then, in the UCRT64 shell:
\begin{lstlisting}[style=shell]
cmake --preset release -DCMAKE_EXE_LINKER_FLAGS=-static && cmake --build --preset release
\end{lstlisting}
(\code{-static}: the program needs no MSYS2 DLLs). The continuous integration builds the C++
route (MSYS2); the GNU Rust toolchain is not tested there.

\paragraph{macOS.} The Command Line Tools need an administrator once.

\section{Without internet access}
\label{sec:installation-offline}

\paragraph{Copy the programs.} Build on a machine with the same operating system and
architecture (on Linux, preferably the static musl build of \cref{sec:installation-rust}) and
copy the single files. Nothing else is needed to run them.

\paragraph{Build on the offline machine.} On a machine with internet access, from the root of
the clone:
\begin{lstlisting}[style=shell]
mkdir -p .cargo && cargo vendor --locked vendor > .cargo/config.toml  # every Rust dependency
tools/fetch-cpp-deps.sh             # the C++ dependencies, into viewers/vitalstream/cpp/deps/
pip download -d ../wheels './viewers/vitalstream/python/vitalstream[wfdb]' "setuptools>=77" wheel
\end{lstlisting}
The last line is needed only for the Python viewer. \cmd{cargo vendor} takes wfdb-x too, and
\file{tools/fetch-cpp-deps.sh} its C++ archive. Copy the whole clone, and \file{../wheels}, to
the offline machine, which needs Rust and a compiler (rustup's offline installers:
\url{https://forge.rust-lang.org/infra/other-installation-methods.html}). Then build as usual:
\cmd{cargo build --release --locked --offline}, \cmd{cmake --preset release}. pip takes wfdb-x
from the archive that \cmd{pip download} saved (named after the commit), then the viewer:
\begin{lstlisting}[style=shell]
WFDBX=a707ba9ec1624241793bf41b938b219d76005526     # the pinned commit of wfdb-x
pip install --no-index --find-links ../wheels fastapi uvicorn numpy \
  "wfdb-x @ file://$PWD/../wheels/$WFDBX.tar.gz#subdirectory=python"
pip install --no-index --find-links ../wheels --no-deps ./viewers/vitalstream/python/vitalstream
\end{lstlisting}
Git ignores \file{vendor/}, \file{.cargo/} and \file{viewers/vitalstream/cpp/deps/}.

\paragraph{The demo, offline.} The installers' \opt{--demo} needs NumPy and wfdb-x in
\file{.venv}; \cmd{pip download} above saved both in \file{../wheels}. Install them from there,
then make and load the demo by hand (\cref{chap:quickstart}):
\begin{lstlisting}[style=shell]
WFDBX=a707ba9ec1624241793bf41b938b219d76005526
.venv/bin/pip install --no-index --find-links ../wheels numpy \
  "wfdb-x @ file://$PWD/../wheels/$WFDBX.tar.gz#subdirectory=python"
.venv/bin/python tools/make-demo-data.py ~/medorama-demo
medorama demo add ~/medorama-demo
\end{lstlisting}

\section{Windows}
\label{sec:installation-windows}

Everything above runs on Windows 10 and 11. In short:
\begin{enumerate}
\item Install Git, Python and the Microsoft C++ Build Tools (\cref{sec:installation-prerequisites}).
\item In PowerShell, from the root of the clone, run
  \cmd{powershell -ExecutionPolicy Bypass -File tools\textbackslash install-medorama.ps1}. Add
  \opt{-Viewer rs} to build the Rust viewer, which needs nothing beyond Rust, and \opt{-Demo} for
  the demo.
\item Activate the environment in each new window:
  \cmd{.venv\textbackslash Scripts\textbackslash Activate.ps1}.
\end{enumerate}
The programs are copied into \file{.venv\textbackslash Scripts}. The C++ viewer is built with
MSVC; built by hand, it needs a Developer PowerShell (\cref{sec:installation-cpp}). The
VitalStream installer for Windows, \file{viewers\textbackslash vitalstream\textbackslash install.ps1},
builds the C++ viewer from a Developer PowerShell for VS. As on Linux and macOS, it builds in
folders of its own (\file{viewers\textbackslash vitalstream\textbackslash cpp\textbackslash build\textbackslash standalone},
\file{target\textbackslash standalone}) and copies the viewer into
\file{\%LOCALAPPDATA\%\textbackslash vitalstream\textbackslash bin}. It has no \opt{-NoRustup}: it never
installs Rust, and \opt{-Viewer rs} fails when \cmd{cargo} is missing.

When PowerShell says ``running scripts is disabled on this system'', start the installers as
shown (\cmd{powershell -ExecutionPolicy Bypass -File \dots}); for \file{Activate.ps1}, allow
local scripts once with \cmd{Set-ExecutionPolicy -Scope CurrentUser RemoteSigned}.

\section{Checking the build}
\label{sec:installation-check}

\begin{lstlisting}[style=shell]
cargo test --release --locked             # every Rust test of the workspace (a few minutes)
(cd viewers/vitalstream/cpp && ctest --preset release)
.venv/bin/python -m pytest -q viewers/vitalstream/python/vitalstream/tests
make parity                               # the three viewers agree (index and answers)
make test-nowfdb                          # the builds without WFDB support
make help                                 # every check
\end{lstlisting}

\section{Updating and uninstalling}
\label{sec:installation-update}

\paragraph{Update.} \cmd{git pull}, then the installer again with the same options (or
\cmd{cargo build --release --locked} and \cmd{cmake --build --preset release}). Without
\opt{--demo}, catalogs created afterwards start empty; a catalog that holds the demo keeps it.
Restart the programs. Medorama updates its database when it starts, and an older version
refuses a database a newer one has updated: back the catalog up first
(\cref{chap:administration}).

\paragraph{Uninstall.} Delete the clone (the programs, \file{.venv/} and the builds are inside
it), or the installed copies (\cmd{cargo uninstall medorama},
\file{\textasciitilde/.local/bin/vitalstream-*}). A standalone viewer's index folder is
\file{\textasciitilde/.cache/vitalstream}, its settings and password are in
\file{\textasciitilde/.config/vitalstream/}; Medorama's home and results folders are where
\cmd{medorama init} put them. Recordings are never touched.

\section{When the build fails}
\label{sec:installation-troubleshooting}

\begin{xltabular}{\linewidth}{@{}P{5.4cm}L@{}}
\caption{Build messages, their causes and fixes.}
\label{tab:installation-troubleshooting}\\
\toprule
{\tabhead Message} & {\tabhead Cause and fix} \\
\midrule
\endfirsthead
\toprule
{\tabhead Message} & {\tabhead Cause and fix} \\
\midrule
\endhead
\bottomrule
\endfoot
\code{command not found: medorama} & The program is in \file{.venv/bin}:
  \cmd{source .venv/bin/activate}, or call \file{.venv/bin/medorama}, or
  \cmd{cargo install --path medorama --locked}. \\
\code{linker 'cc' not found} (Linux) & No C toolchain: \code{build-essential} or \code{gcc}
  (\cref{sec:installation-prerequisites}), or \cref{sec:installation-no-admin}. \\
\code{linker 'link.exe' not found}, \code{No CMAKE\_CXX\_COMPILER could be found} (Windows) &
  The Build Tools are missing or lack \emph{Desktop development with C++}; for C++, use a
  Developer PowerShell. \\
\code{exit status: 69}, \code{You have not agreed to the Xcode license agreements} (macOS) &
  \cmd{sudo xcodebuild -license accept}, or
  \cmd{export DEVELOPER\_DIR=/Library/Developer/CommandLineTools}. \\
\code{feature 'edition2024' is required}, \code{requires rustc 1.88} & Rust is too old:
  \cmd{rustup update} (a distribution's packaged Rust is often old). \\
\code{cargo: command not found} after installing Rust & Open a new shell, or
  \cmd{. \textasciitilde/.cargo/env}. \\
\code{failed to get \dots{} from crates.io} & No internet access: vendor the dependencies
  (\cref{sec:installation-offline}). \\
\code{The C++ compiler is too old} & GCC 12+ or Clang 15+ is needed
  (\cref{sec:installation-cpp}, \cref{sec:installation-no-admin}); then delete the build
  folder. \\
\code{CXX is \dots, but this build folder was configured with \dots} & CMake keeps the first
  compiler: \cmd{rm -rf build/release} and configure again. \\
\code{The current CMakeCache.txt directory \dots{} is different} & The clone moved after a
  build: delete \file{build/} and configure again. \\
\code{CMake 3.24 or higher is required}, no \code{Ninja} &
  \cmd{pip install --user "cmake>=3.24" ninja}. \\
\code{Failed to download}, \code{HASH mismatch} while configuring & No internet access or a
  proxy: \file{tools/fetch-cpp-deps.sh} on another machine (\cref{sec:installation-offline}). \\
\code{version 'GLIBC\_2.xx' not found}, \code{'GLIBCXX\_3.4.xx' not found} for a copied program &
  Built on a newer Linux: the static musl build (Rust), or the static C++ runtime (C++). \\
``running scripts is disabled on this system'' (Windows) & Start the installers as shown
  (\cmd{powershell -ExecutionPolicy Bypass -File \dots}); for \file{Activate.ps1}:
  \cmd{Set-ExecutionPolicy -Scope CurrentUser RemoteSigned}. \\
\end{xltabular}

Problems that appear after the build, when running a catalog, are covered in
\cref{chap:troubleshooting}.
